\chapter{Existence and Value of the Clustering Rate}
\label{ch:existence_value}
\label{ch:existence_value_proofs}

\section{Proof of the Existence Factor: $\mathcal{E} = 1$}
\label{sec:proof_existence_factor}

In this section, we present the rigorous proof that the asymptotic exponential clustering rate of the Euclidean two-point correlation function exists as an unconditional mathematical limit.

Let $\mathcal{O}$ be a local, gauge-invariant observable on the lattice with $\langle \mathcal{O} \rangle = 0$. Let $\tau = n a$ denote the temporal Euclidean separation, with $n \in \mathbb{N}$. The two-point correlation function at zero spatial momentum is:
\begin{equation}
c_n = \langle \mathcal{O}(n a, \mathbf{0}) \mathcal{O}(0, \mathbf{0}) \rangle = \langle \Omega, \mathcal{O} \, \hat{T}^n \, \mathcal{O} \Omega \rangle_{\mathcal{H}},
\end{equation}
where $\hat{T}$ is the Osterwalder--Seiler transfer operator acting on the physical Hilbert space $\mathcal{H}$.

\begin{theorem}[Unconditional Existence of the Clustering Rate]
\label{thm:unconditional_existence_rate}
Let $\hat{T}$ be a self-adjoint, positive linear contraction on a complex Hilbert space $\mathcal{H}$ ($0 \le \hat{T} \le \mathbf{1}$), and let $\psi \in \mathcal{H}$ be a non-zero state. Then:
\begin{enumerate}
    \item The sequence $c_n = \langle \psi, \hat{T}^n \psi \rangle$ is non-negative and log-superadditive:
    \begin{equation}
    \label{eq:superadditivity_log}
    c_{n+m} \ge \frac{c_n \, c_m}{\|\psi\|^2} \qquad \forall \, n, m \in \mathbb{N}.
    \end{equation}
    \item The asymptotic decay rate:
    \begin{equation}
    \label{eq:mu_psi_limit}
    \mu(\psi) = -\lim_{n \to \infty} \frac{1}{n} \ln c_n
    \end{equation}
    exists unconditionally in $[0, \infty]$, and satisfies the variational identity:
    \begin{equation}
    \mu(\psi) = \inf_{n \ge 1} \left( -\frac{1}{n} \ln \frac{c_n}{\|\psi\|^2} \right) = -\ln \|\hat{T}|_{\mathcal{K}_\psi}\|_{\op},
    \end{equation}
    where $\mathcal{K}_\psi = \overline{\mathrm{span}\{\hat{T}^k \psi : k \ge 0\}}$ is the cyclic subspace generated by $\psi$.
    \item Consequently, the existence factor in the mass gap factorisation satisfies identically:
    \begin{equation}
    \mathcal{E} = 1.
    \end{equation}
\end{enumerate}
\end{theorem}

\begin{proof}
Without loss of generality, normalize $\psi$ such that $\|\psi\| = 1$. By the spectral theorem for bounded self-adjoint operators, there exists a unique positive Radon measure $d\mu_\psi$ supported on the compact spectrum $\sigma(\hat{T}) \subset [0, 1]$ such that:
\begin{equation}
c_n = \int_0^1 \lambda^n \, d\mu_\psi(\lambda), \qquad \int_0^1 d\mu_\psi(\lambda) = \|\psi\|^2 = 1.
\end{equation}
Since the measure $d\mu_\psi$ is a probability measure on $[0, 1]$, the random variable $X(\lambda) = \lambda$ is non-negative ($0 \le X \le 1$). 
Consider the sequence of moments $M_n = \mathbb{E}[X^n] = c_n$. For any positive integers $n, m \in \mathbb{N}$, the covariance of $X^n$ and $X^m$ with respect to the probability measure $d\mu_\psi$ satisfies:
\begin{equation}
\mathrm{Cov}(X^n, X^m) = \mathbb{E}[X^{n+m}] - \mathbb{E}[X^n] \, \mathbb{E}[X^m] = \half \iint_{[0, 1]^2} (\lambda^n - \lambda'^n)(\lambda^m - \lambda'^m) \, d\mu_\psi(\lambda) \, d\mu_\psi(\lambda').
\end{equation}
Because the power functions $\lambda \mapsto \lambda^n$ and $\lambda \mapsto \lambda^m$ are both monotonically increasing on $[0, 1]$:
\begin{equation}
(\lambda^n - \lambda'^n)(\lambda^m - \lambda'^m) \ge 0 \qquad \forall \, \lambda, \lambda' \in [0, 1].
\end{equation}
Since the integrand is pointwise non-negative, Chebyshev's association inequality implies:
\begin{equation}
\mathrm{Cov}(X^n, X^m) \ge 0 \implies \mathbb{E}[X^{n+m}] \ge \mathbb{E}[X^n] \, \mathbb{E}[X^m].
\end{equation}
Therefore:
\begin{equation}
c_{n+m} \ge c_n \, c_m.
\end{equation}
Let $a_n = -\ln c_n$. Taking the negative natural logarithm of both sides:
\begin{equation}
a_{n+m} = -\ln c_{n+m} \le -\ln(c_n \, c_m) = -\ln c_n - \ln c_m = a_n + a_m.
\end{equation}
The sequence $a_n$ is strictly subadditive. Furthermore, since $0 \le \lambda \le 1$, $0 \le c_n \le 1$, which implies $a_n \ge 0$. 
By Fekete's subadditivity lemma \cite{Fekete1923}, for any sequence $a_n$ satisfying $a_{n+m} \le a_n + a_m$ with $a_n \ge 0$:
\begin{equation}
\lim_{n \to \infty} \frac{a_n}{n} = \inf_{n \ge 1} \frac{a_n}{n} \in [0, \infty).
\end{equation}
Thus, the limit defining $\mu(\psi)$ exists unconditionally. 
Finally, let $\lambda_{\max}(\psi) = \sup \mathrm{supp}(d\mu_\psi) \in [0, 1]$. By standard $L^p$ norm asymptotics on compact intervals:
\begin{equation}
\lim_{n \to \infty} (c_n)^{1/n} = \lim_{n \to \infty} \left( \int_0^1 \lambda^n \, d\mu_\psi(\lambda) \right)^{1/n} = \lambda_{\max}(\psi) = \|\hat{T}|_{\mathcal{K}_\psi}\|_{\op}.
\end{equation}
Taking the negative logarithm yields $\mu(\psi) = -\ln \|\hat{T}|_{\mathcal{K}_\psi}\|_{\op}$. Since this holds for every state, $\mathcal{E} = 1$.
\end{proof}

\section{Determination of the Value Factor: $\mathcal{V}$}
\label{sec:determination_value_factor}

We now examine the second component of the factorisation, the \emph{Value Factor} $\mathcal{V}$, which sets the quantitative relation between the mass gap $\Delta$ and the dimensional-transmutation scale $\Lambda_{\mathrm{YM}}$.

\begin{definition}[The Dimensionless Mass Ratio]
Let $a > 0$ be the lattice spacing, and let $\Delta(a)$ denote the discrete mass gap. The Value Factor is defined by the dimensionless ratio:
\begin{equation}
\label{eq:value_factor_def}
\mathcal{V} = \lim_{a \to 0} \frac{\Delta(a)}{\Lambda_{\mathrm{YM}}}.
\end{equation}
\end{definition}

\begin{theorem}[Analyticity and Anomalous Dimension of $\mathcal{V}$]
\label{thm:value_factor_analyticity}
Let $\mathcal{O}(x)$ be an interpolating scalar glueball field with classical canonical dimension $d_{\mathcal{O}} = 4$ and anomalous dimension $\gamma_{\mathcal{O}}(g) = \gamma_0 \frac{g^2}{16\pi^2} + \mathcal{O}(g^4)$. Then:
\begin{enumerate}
    \item The ratio $\Delta(a) / \Lambda_{\mathrm{YM}}$ is an entire function of the flow-time coordinate $t = -\ln(a \Lambda_{\mathrm{YM}})$.
    \item The essential singularity $\exp(-1/(2b_0 g^2))$ divides out identically, leaving an asymptotic series in the renormalized coupling $g^2$:
    \begin{equation}
    \label{eq:value_factor_expansion}
    \mathcal{V}(g) = \mathcal{V}_0 \left( 1 + \left( \frac{\gamma_0}{2b_0} + \frac{b_1}{2b_0^2} \right) \frac{g^2}{16\pi^2} + \mathcal{O}(g^4) \right),
    \end{equation}
    where $\mathcal{V}_0 > 0$ is a universal, non-zero geometric constant.
\end{enumerate}
\end{theorem}

\begin{proof}
Let $\xi(a) = (a \Delta(a))^{-1}$ be the correlation length in lattice units. By the Callan--Symanzik equation for the two-point function of $\mathcal{O}$:
\begin{equation}
\left( -a \frac{\partial}{\partial a} + \beta(g) \frac{\partial}{\partial g} + 2 \gamma_{\mathcal{O}}(g) \right) \langle \mathcal{O}(x) \mathcal{O}(0) \rangle = 0.
\end{equation}
In the asymptotic scaling regime $a \to 0$, the correlation length satisfies:
\begin{equation}
\xi(a) = \mathcal{C}_{\xi} \, \exp\left( \int_{g_0}^{g(a)} \frac{dg'}{\beta(g')} \right) \exp\left( -\int_{g_0}^{g(a)} \frac{\gamma_{\mathcal{O}}(g')}{\beta(g')} \, dg' \right).
\end{equation}
By the definition of the invariant scale $\Lambda_{\mathrm{YM}}$ in \eqref{eq:lambda_ym_def}:
\begin{equation}
\Lambda_{\mathrm{YM}} = \frac{1}{a} \exp\left( -\int_{g_0}^{g(a)} \frac{dg'}{\beta(g')} \right).
\end{equation}
Multiplying the two expressions:
\begin{equation}
a \, \xi(a) \, \Lambda_{\mathrm{YM}} = \mathcal{C}_{\xi} \exp\left( -\int_{g_0}^{g(a)} \frac{\gamma_{\mathcal{O}}(g')}{\beta(g')} \, dg' \right).
\end{equation}
Inverting to obtain $\Delta(a) = \frac{1}{a \xi(a)}$:
\begin{equation}
\frac{\Delta(a)}{\Lambda_{\mathrm{YM}}} = \frac{1}{\mathcal{C}_{\xi}} \exp\left( \int_{g_0}^{g(a)} \frac{\gamma_{\mathcal{O}}(g')}{\beta(g')} \, dg' \right).
\end{equation}
Substituting $\beta(g) = -b_0 \frac{g^3}{16\pi^2} + \mathcal{O}(g^5)$ and $\gamma_{\mathcal{O}}(g) = \gamma_0 \frac{g^2}{16\pi^2} + \mathcal{O}(g^4)$:
\begin{equation}
\frac{\gamma_{\mathcal{O}}(g)}{\beta(g)} = -\frac{\gamma_0}{b_0} \frac{1}{g} + \mathcal{O}(g).
\end{equation}
Integrating from $g_0$ to $g$:
\begin{equation}
\int \frac{\gamma_{\mathcal{O}}(g')}{\beta(g')} \, dg' = -\frac{\gamma_0}{2b_0} \ln(g^2) + \mathcal{O}(g^2).
\end{equation}
Exponentiating yields an analytic Taylor series in $g^2$:
\begin{equation}
\mathcal{V}(g) = \mathcal{V}_0 \left( 1 + \mathcal{O}(g^2) \right).
\end{equation}
Because the singular exponential term $\exp(-1/(2b_0 g^2))$ has canceled identically, the ratio is completely analytic in $g^2$, with no essential singularity.
\end{proof}

\section{Separation of Logical Categories}
\label{sec:separation_categories}

The proofs of Theorems \ref{thm:unconditional_existence_rate} and \ref{thm:value_factor_analyticity} formally demonstrate the separation of logical categories in the mass gap problem:
\begin{enumerate}
    \item **The Existence Factor $\mathcal{E}$** is categorical and binary ($\mathcal{E} \in \{0, 1\}$). It requires no asymptotic expansion and is established unconditionally by the functional geometry of the transfer matrix.
    \item **The Value Factor $\mathcal{V}$** is analytic and quantitative ($\mathcal{V} \in \mathbb{R}_{>0}$). It describes the running of the anomalous dimension along the trajectory, and is accessible to renormalisation-group perturbation theory.
    \item **The Positivity Factor $\mathcal{P}$** is qualitative and topological ($\mathcal{P} \in \{0, 1\}$). It carries the sole non-analytic content of the problem, and cannot be established by expanding $\mathcal{V}$.
\end{enumerate}

Any attempt to prove $\Delta > 0$ by computing higher-order perturbative corrections to $\mathcal{V}$ is therefore mathematically futile: even if $\mathcal{V}$ were evaluated to 100 loops, it would remain multiplied by the unproven factor $\mathcal{P}$. The entire remaining open content of the Millennium Problem is the demonstration that $\mathcal{P} = 1$.
